\scanpage{8}
à prouver que si $f:X\mathbin{-}Y$ fait de $\unsure{X|x}$ un rev. fini étale de $Y$ et si $u\in\mathcal{A}_X^0$ est tel que $f=f\circ u$ i.e. $u$ est un $Y$-morphisme, alors $u=\operatorname{id}$. Mais c'est ce qu'affirme la prop. 2.

\uline{Corollaire 2} Les composantes connexes de l'espace des morphismes finis étales $X\longrightarrow Y$ sont contractiles.

\uline{Corollaire 3} \deleted{L'ensemble $\operatorname{Hom}_{\mathrm{isot}}(X,Y)$ est un torseur \inserted{\deleted{à droite}} sous le} \inserted{L'opération à droite du} groupe de Teichmüller (non orienté) $\unsure{\mathcal{T}_X|T_X}=\mathcal{A}_X/\mathcal{A}_X^0$ sur l'ens $\operatorname{Hom}_{\mathrm{isot}}(X,Y)$ des ``morphismes isotopiques'' de $X$ dans $Y$ \deleted{est libre,} a des stabilisateurs qui sont des sous-groupes finis de $\mathcal{T}_X$.
\margin{\uline{NB} Tout sous-groupe fini de $\mathcal{T}_X$ peut s'obtenir de cette façon comme stabilisateur !! (th. profond et récent\ldots)}

De façon précise, le stabilisateur de $f\in\operatorname{Hom}_{\mathrm{isot}}(X,Y)$ s'identifie au \deleted{\unsure{un}} \inserted{groupe} \deleted{\unsure{gr.}} des $Y$-autom. du rev. fini étale $Y'_f$ défini par $f$. L'ensemble des orbites de $\mathcal{T}_X$ opérant sur $\operatorname{Hom}_{\mathrm{isot}}(X,Y)$ est en corr. biunivoque avec l'ens \inserted{$\deleted{\unsure{\operatorname{Rev}(Y)}}\ \unsure{R_X(Y)|\mathcal{R}_X(Y)}$} des classes d'isom. de revêtements finis étales de $Y$, homéomorphes à $X$.
\margin{\uline{NB} cet ens d'orbites est \uuline{fini}}
\figure{\begin{tikzpicture}[baseline=0pt]\draw (0,0)--(0,2.1);\end{tikzpicture}}
[L'ens des orbites sous $\mathcal{T}_X^+$ est en corr. biunivoque avec l'ens. $R_X(Y)\times\{\pm1\}$, le signe $+$ \unsure{resp. $-$} correspondant aux classes de morphismes finis étales $X\longrightarrow Y$
\margin{7}
